\documentclass[11pt]{article}
\usepackage{amsmath,amssymb,amsthm,mathtools}
\usepackage[margin=1in]{geometry}
\usepackage{enumitem}
\usepackage{booktabs}
\usepackage{hyperref}

\newtheorem{theorem}{Theorem}[section]
\newtheorem{proposition}[theorem]{Proposition}
\newtheorem{lemma}[theorem]{Lemma}
\newtheorem{definition}[theorem]{Definition}
\newtheorem{corollary}[theorem]{Corollary}
\newtheorem{example}[theorem]{Example}
\newtheorem{warning}[theorem]{Warning}

\providecommand{\Cap}{\operatorname{Cap}}
\newcommand{\Ran}{\operatorname{Ran}}
\newcommand{\Tr}{\operatorname{Tr}}
\newcommand{\spec}{\operatorname{spec}}
\newcommand{\diag}{\operatorname{diag}}
\newcommand{\End}{\operatorname{End}}
\newcommand{\Ker}{\operatorname{Ker}}
\newcommand{\id}{\operatorname{id}}
\newcommand{\C}{\mathbb C}
\newcommand{\K}{\mathsf K}
\newcommand{\D}{\mathsf D}
\newcommand{\G}{\mathcal G}
\newcommand{\cH}{\mathcal H}

\title{Step 31: Root-Composite / Algebraic-Extension Anti-Localization Profile}
\author{Six Birds Anti-Localization Project}
\date{May 2026}

\begin{document}
\maketitle

\section{Purpose}
The scalar and family currency theorems developed so far are analytic: probes are represented by linear functionals, and capacity is the shadow price matrix
\[
  \K=L_\Gamma C_\Gamma^\dagger L_\Gamma^* .
\]
The corpus reading identifies a missing fourth profile of anti-localization: the \emph{root-composite} or algebraic-extension profile.  This note formulates that profile.  It is not an RH proof.  It is a theorem schema for the case in which the anti-localization carrier lives in an algebraic extension and lower-layer claims descend through traces, characters, norms, or characteristic-polynomial data.

The guiding idea is:
\[
\boxed{\text{root-composite anti-localization is character-wise pricing of an algebraic orbit of probes.}}
\]

\section{Algebraic extension carriers}

\begin{definition}[finite algebraic extension carrier]
A finite algebraic anti-localization carrier is a tuple
\[
  \mathfrak A=(H,C,\Gamma,G,U,\rho,L,\Theta),
\]
where:
\begin{enumerate}[label=(\roman*)]
  \item $H$ is a finite-dimensional Hilbert carrier;
  \item $C\succeq0$ is the audit energy;
  \item $\Gamma:E\to H$ is the exact package map and $C_\Gamma=\Gamma^*C\Gamma$;
  \item $G$ is a finite group acting unitarily on $H$ by $U_g$;
  \item $\rho:G\to U(Y)$ is a finite-dimensional unitary response representation;
  \item $L:H\to Y$ is a declared probe family satisfying the equivariance law
  \[
    L U_g=\rho_g L;
  \]
  \item $\Theta\succeq0$ is a declared response budget, normally required to commute with $\rho(G)$.
\end{enumerate}
The package is algebraically lawful when the audit and package are compatible with the group action:
\[
  U_g^*CU_g=C,
\]
and, when the package coordinates are used, $\Ran\Gamma$ is $G$-stable or has a declared transported quotient action.
\end{definition}

\begin{definition}[root-composite currency matrix]
The root-composite currency matrix of an algebraic extension carrier is
\[
  \K_{G}=L_\Gamma C_\Gamma^\dagger L_\Gamma^*,
  \qquad L_\Gamma=L\Gamma.
\]
For $y\in Y$, the recombined probe is
\[
  \ell_y(u)=\langle y,Lu\rangle_Y.
\]
The response $y$ may be a trace mode, a character mode, or a more general algebraic recombination.
\end{definition}

\section{Equivariant descent and character blocks}

\begin{theorem}[equivariant currency descent]
Assume $C_\Gamma\succeq0$, $L_\Gamma$ annihilates $\ker C_\Gamma$, and the package/audit/probe data are $G$-equivariant.  Then
\[
  \rho_g \K_G \rho_g^*=\K_G
  \qquad(g\in G).
\]
Consequently, $\K_G$ belongs to the commutant of the response representation $\rho$ and decomposes into isotypic blocks.  If
\[
  Y=\bigoplus_{\pi\in\widehat G}Y_\pi
\]
is the isotypic decomposition, then
\[
  \K_G=\bigoplus_{\pi\in\widehat G}\K_\pi
\]
up to multiplicity blocks in the commutant.
\end{theorem}

\begin{proof}
Equivariance gives $L_\Gamma T_g=\rho_gL_\Gamma$ for the induced action $T_g$ on package coordinates.  Audit equivariance gives $T_g^*C_\Gamma T_g=C_\Gamma$ and hence $T_g^*C_\Gamma^\dagger T_g=C_\Gamma^\dagger$.  Therefore
\[
  \rho_g\K_G\rho_g^*
  =\rho_g L_\Gamma C_\Gamma^\dagger L_\Gamma^*\rho_g^*
  =L_\Gamma T_g C_\Gamma^\dagger T_g^* L_\Gamma^*
  =L_\Gamma C_\Gamma^\dagger L_\Gamma^*
  =\K_G .
\]
The block decomposition follows from complete reducibility of finite-dimensional unitary representations and Schur's lemma applied to the commutant.
\end{proof}

\begin{corollary}[abelian character ledger]
If $G$ is abelian and the response representation is a direct sum of characters, then the root-composite budget condition
\[
  \K_G\preceq\Theta
\]
is equivalent to the character-wise inequalities
\[
  \K_\chi\preceq\Theta_\chi
  \qquad(\chi\in\widehat G).
\]
If $Y$ is the regular representation of a cyclic group $C_m$ and $\Theta=I$, then $\K_G$ is circulant and
\[
  \K_G\preceq I
  \iff
  \widehat k(\ell)\le1\quad(\ell=0,\dots,m-1),
\]
where $\widehat k$ is the discrete Fourier transform of the first row of $\K_G$.
\end{corollary}

\section{Capacity, trace modes, and character witnesses}

\begin{theorem}[root-composite capacity identity]
For every algebraic recombination $y\in Y$,
\[
  \Cap_C(\ell_y;\Ran\Gamma)=y^*\K_Gy.
\]
Thus the extension carrier is anti-localized against all algebraic recombinations exactly when
\[
  \K_G\preceq\Theta.
\]
If the inequality fails, then there is a character/isotypic witness $y_\pi\in Y_\pi$ such that
\[
  y_\pi^*(\K_G-\Theta)y_\pi>0.
\]
\end{theorem}

\begin{proof}
The first identity is the family currency theorem applied to $L_\Gamma$.  If $\K_G\npreceq\Theta$, then $\D=\K_G-\Theta$ has a positive Rayleigh quotient.  When $\Theta$ is equivariant, $\D$ is equivariant and decomposes by isotypic blocks; a positive eigenvector may therefore be chosen in one block or in a multiplicity component inside that block.
\end{proof}

\begin{definition}[trace mode and character modes]
For an orbit probe family indexed by $G$, the normalized trace response is
\[
  y_{\mathrm{tr}}=|G|^{-1/2}\sum_{g\in G}e_g.
\]
It is the trivial character mode.  Nontrivial character modes are the remaining components of the regular representation.  Testing only $y_{\mathrm{tr}}$ is a trace-shadow check; testing all characters is the root-composite check.
\end{definition}

\begin{warning}[trace descent is not enough]
The trace mode can pass while a nontrivial character mode fails.  Let
\[
  \K=(1-c)I+cJ_m,
  \Theta=I,
\]
where $J_m$ is the all-ones matrix.  The coordinate capacities are all $1$.  The trivial character eigenvalue is
\[
  \lambda_{\mathrm{tr}}=1+(m-1)c,
\]
and each nontrivial character eigenvalue is
\[
  \lambda_{\mathrm{nt}}=1-c.
\]
For $-1/(m-1)<c<0$, the matrix is positive semidefinite, the trace mode has $\lambda_{\mathrm{tr}}<1$, but every nontrivial character mode has $\lambda_{\mathrm{nt}}>1$.  Thus trace descent can miss a root-composite needle.
\end{warning}

\section{Characteristic-polynomial descent}

\begin{definition}[root ledger]
Assume $\Theta\succ0$.  The relative root-composite operator is
\[
  R=\Theta^{-1/2}\K_G\Theta^{-1/2}.
\]
Its characteristic polynomial
\[
  p_R(z)=\det(zI-R)
\]
records the descended root ledger.  The accepted anti-localization condition is
\[
  \lambda_{\max}(R)\le1.
\]
When $G$ is abelian, this root ledger factors into character polynomials.
\end{definition}

\begin{theorem}[root ledger criterion]
For $\Theta\succ0$,
\[
  \K_G\preceq\Theta
  \iff
  \lambda\le1\quad\text{for every root }\lambda\text{ of }p_R.
\]
Equivalently, every character/block root ledger lies in $(-\infty,1]$.
\end{theorem}

\begin{proof}
Since $R$ is self-adjoint positive semidefinite, $\K_G\preceq\Theta$ is equivalent to $R\preceq I$, which is equivalent to all eigenvalues of $R$ being at most $1$.  These eigenvalues are the roots of $p_R$ counted with multiplicity.  If the data are equivariant, $R$ decomposes into isotypic blocks and the characteristic polynomial factors accordingly.
\end{proof}

\begin{warning}[norm or determinant shadows are insufficient]
A determinant or norm descent can be a public shadow rather than a lawful root ledger.  For example,
\[
  R=\begin{pmatrix}2&0\\0&1/4\end{pmatrix}
\]
has determinant $1/2<1$, but $R\npreceq I$ because its largest root is $2$.  Root-composite anti-localization requires a largest-root or Loewner certificate, not merely a product norm.
\end{warning}

\section{Algebraic descent as a strict-extension profile}

The root-composite profile differs from the invariant, coefficient, and obstruction-ledger profiles as follows.
\begin{center}
\begin{tabular}{lll}
\toprule
Profile & Carrier lives in & Descent expression \\
\midrule
Invariant & extended operator family & trace / sector capacity \\
Coefficient & infinitesimal or perturbative extension & standard part / coefficient \\
Obstruction-ledger & cohomological or route extension & obstruction class / union ledger \\
Root-composite & algebraic/Galois extension & roots / characters / characteristic ledger \\
\bottomrule
\end{tabular}
\end{center}

The root-composite profile is therefore not ``more Fourier analysis.''  It is a strict-extension mode in which the native probes are algebraic conjugates and the lower-layer anti-localization statement is obtained by a full root ledger.

\section{Closure-only theorem}

\begin{theorem}[root-composite no-native-root-needle theorem]
Let $\mathfrak A=(H,C,\Gamma,G,U,\rho,L,\Theta)$ be a formed algebraic closure: exact package, declared group action, equivariant audit, native orbit probe family, null-mode legality, and accepted response budget.  If
\[
  \K_G=L_\Gamma C_\Gamma^\dagger L_\Gamma^*\preceq\Theta,
\]
then no native algebraic recombination has over-budget capacity.  Equivalently, there is no native root-composite critical-pair witness.
\end{theorem}

\begin{proof}
For every native algebraic recombination $y$, the capacity identity gives
\[
  \Cap_C(\ell_y;\Ran\Gamma)=y^*\K_Gy\le y^*\Theta y.
\]
Thus no $y$ can be an over-budget witness.  Conversely, if such a witness existed, $\K_G\npreceq\Theta$ by the separating-vector criterion.
\end{proof}

\section{Audit statuses and nonclaims}

A root-composite anti-localization record must include:
\begin{enumerate}[label=(\roman*)]
  \item the extension law $G$ or algebraic extension $B/A$;
  \item the action on the carrier and on the response family;
  \item audit equivariance $U_g^*CU_g=C$;
  \item package equivariance or a transported quotient action;
  \item the full character/root budget, not only a trace or determinant shadow;
  \item a nonclaim forbidding literalization: a root-composite bound licenses no native algebraic recombination needle under this audit, not the statement that the substrate has no abstract needle-like eigenvectors.
\end{enumerate}

\section{Numerical sanity checks}
The accompanying script checks four finite identities and countermodels:
\begin{enumerate}
  \item cyclic character modes can fail although all coordinate probes have unit scalar capacity;
  \item circulant currency matrices diagonalize by characters;
  \item determinant shadows can pass while Loewner budgets fail;
  \item equivariance defects block algebraic descent.
\end{enumerate}
These checks are not Six Birds simulations.  They are finite sanity checks for the algebraic identities above.

\end{document}
